\documentclass[12pt]{article}
\usepackage{amsfonts,amsthm,amsmath,amssymb,mathtools}
\usepackage{mathrsfs}
\usepackage{enumitem}
\usepackage{tikz-cd}
\usepackage{geometry}
\usepackage{graphicx}
\usepackage{mlmodern}

\geometry{letterpaper, portrait, margin=1in}

\setcounter{section}{0}

\renewcommand{\thesection}{\S\arabic{section}}
\renewcommand{\thesubsection}{\arabic{section}}
\DeclareMathOperator{\im}{im}
\DeclareMathOperator{\id}{id}
\DeclareMathOperator{\ad}{ad}
\DeclareMathOperator{\tr}{tr}
\newcommand{\gl}{\mathfrak{gl}}
\newcommand{\so}{\mathfrak{so}}
\renewcommand{\sl}{\mathfrak{sl}}
\renewcommand{\th}{^\mathrm{th}}

\newtheorem{thm}{Theorem}
\newenvironment{theorem}[1]{
	\renewcommand{\thethm}{#1}
	\begin{thm}
		}{\end{thm}}

\newtheorem{lma}{Lemma}
\newenvironment{lemma}[1]{
	\renewcommand{\thelma}{#1}
	\begin{lma}
		}{\end{lma}}

\theoremstyle{definition}
\newtheorem{ex}{}
\newenvironment{exercise}[1]{
	\renewcommand{\theex}{#1}
	\begin{ex}
		}{\end{ex}}


\title{Lie Theory: Chapter 8}
\author{Connor Mooneyhan}
\date{}

\begin{document}

\maketitle

\begin{ex}
	Consider the usual basis for $\sl(2, \mathbb{C})$: \begin{equation*}
		e = \begin{pmatrix} 0 & 1 \\ 0 & 0 \end{pmatrix}, \hspace{2em} f = \begin{pmatrix} 0 & 0 \\ 1 & 0 \end{pmatrix}, \hspace{2em} h = \begin{pmatrix} 1 & 0 \\ 0 & -1 \end{pmatrix}.
	\end{equation*}
	Fix $t$ (which you can think of as a real number).
	\begin{enumerate}[label=(\alph*)]
		\item Compute $\exp(th)$ and determine what it does to vectors in $\mathbb{R}^2$.
		\item Do the same for $\exp(te)$ and $\exp(tf)$.
	\end{enumerate}
\end{ex}
\begin{proof}\
	\begin{enumerate}[label=(\alph*)]
		\item
		\item
	\end{enumerate}
\end{proof}

\begin{ex}
	If $\mathbb{C}[x_1, \dots, x_n]$ is the vector space of polynomials in $n$ variables and $V_d$ is the subspace of homogeneous polynomials of degree $d$, what is the dimension of $V_d$ as a function of $n$ and $d$?
\end{ex}
\begin{proof}

\end{proof}

\begin{ex}
	Let $\varphi : \sl(2, \mathbb{C}) \to \gl(V_d)$ be the representation of Theorem 8.1.
	Check that $[\varphi(h), \varphi(f)] = -2\varphi(f)$.
\end{ex}
\begin{proof}

\end{proof}

\begin{ex}
	Check that the $\sl(2, \mathbb{C})$-submodule of $V_d$ generated by any particular basis element $X^aY^b$ contains all the basis elements and so is all of $V_d$.
\end{ex}
\begin{proof}

\end{proof}

\begin{ex}
	Find explicit isomorphisms between
	\begin{enumerate}[label=(\roman*)]
		\item the trivial representation of $\sl(2, \mathbb{C})$ and $V_0$;
		\item the natural representation of $\sl(2, \mathbb{C})$ and $V_1$;
		\item the adjoint representation of $\sl(2, \mathbb{C})$ and $V_2$.
	\end{enumerate}
	(Show how you identified the isomorphisms but you don't have to prove they are isomorphisms.)
\end{ex}
\begin{proof}\
	\begin{enumerate}[label=(\roman*)]
		\item
		\item
		\item
	\end{enumerate}
\end{proof}

\begin{ex}
	Exercise 8.4 asserts that if $V$ is a finite-dimensional $\sl(2, \mathbb{C})$-module which is the direct sum of $k$ irreducible modules, then \begin{equation*}
		k = \dim W_0 + \dim W_1,
	\end{equation*}
	where $W_r = \lbrace v \in V : h \cdot v = rv \rbrace$.
	Check this assertion for $V = \sl(3, \mathbb{C})$.
	That is, check that when the finite-dimensional $\sl(2, \mathbb{C})$-module is $\sl(3, \mathbb{C})$, the sum of the dimensions of $W_0$ and $W_1$ equals the number of summands in the decomposition into a sum of irreducibles.
\end{ex}
\begin{proof}

\end{proof}

\begin{ex}
	Let $V$ be an $\sl(2, \mathbb{C})$-module, not necessarily finite-dimensional.
	Suppose $w \in V$ is a highest-weight vector of weight $\lambda$; that is, $e \cdot w = 0$ and $h \cdot w = \lambda w$ for some $\lambda \in \mathbb{C}$, and $w \neq 0$.
	Show that
	\begin{enumerate}[label=(\roman*)]
		\item for $k = 1, 2, \dots$ we have $e \cdot (f^k \cdot w) = k(\lambda - k + 1)f^{k-1} \cdot w$, and
		\item $e^kf^k \cdot w = (k!)^2\binom{\lambda}{k}w$.
	\end{enumerate}
	Deduce that if $\binom{\lambda}{k} \neq 0$ then the set of all $f^j \cdot w$ for $0 \leq j \leq k$ is linearly independent.
	Hence show that if $V$ is finite-dimensional, then $\lambda$ must be a non-negative integer.
\end{ex}
\begin{proof}\
	\begin{enumerate}[label=(\roman*)]
		\item
		\item
	\end{enumerate}
\end{proof}

\end{document}
